% Generated draft: requires mathematical and source review.
\documentclass[11pt]{article}
\usepackage[margin=25mm]{geometry}
\usepackage{fontspec}
\setmainfont{Latin Modern Roman}
\usepackage{amsmath,amssymb,mathtools,unicode-math}
\setmathfont{Latin Modern Math}
\usepackage{enumitem,needspace,xurl,hyperref}
\setlength{\parindent}{0pt}
\setlength{\parskip}{5pt}
\setlength{\emergencystretch}{3em}
\newcommand{\N}{\mathbb{N}}
\newcommand{\Z}{\mathbb{Z}}
\newcommand{\Q}{\mathbb{Q}}
\newcommand{\R}{\mathbb{R}}
\newcommand{\C}{\mathbb{C}}
\newcommand{\F}{\mathbb{F}}
\newcommand{\A}{\mathbb{A}}
\newcommand{\PP}{\mathbb{P}}
\newcommand{\E}{\mathbb{E}}
\newcommand{\eps}{\varepsilon}
\newcommand{\dd}{\,\mathrm{d}}
\begin{document}
\section*{UnsolvedMath ID: 1900021}
\textbf{Source problem code:} catalogue entry
\subsection{Definitions and mathematical statement}
Under the classical normalization, let $Q(x,y)=ax^2+bxy+cy^2$ be a real indefinite binary quadratic form with discriminant $b^2-4ac=1$, and define
\[
 m(Q)=\inf\{|Q(u,v)|:(u,v)\in\Z^2\setminus\{(0,0)\}\}.
\]
The finite Markov spectrum is
\[
 M=\{m(Q)^{-1}:b^2-4ac=1,\ m(Q)>0\}\subset\R.
\]
Forms with $m(Q)=0$ can instead be represented by an infinite spectral value if an extended spectrum is intended. The source does not specify its normalization or whether that value is included.

For the supplied comparison with the Lagrange spectrum, a classical finite-value convention is
\[
 L=\left\{\limsup_{q\to\infty}\frac{1}{q\|q\alpha\|}:\alpha\in\R\setminus\Q,\ 
 \limsup_{q\to\infty}\frac{1}{q\|q\alpha\|}<\infty\right\},
\]
where $q$ runs through positive integers and $\|t\|=\inf_{z\in\Z}|t-z|$.

\textbf{Research direction.} Study the geometric properties of the Markov spectrum. With the above convention, this concerns the geometry of the subset $M\subset\R$, but the source supplies no selected property, conjectural formula, parameter range, or criterion for a complete answer. It is therefore an open-ended program rather than a single quantified assertion. The research summary identifies fractal/Hausdorff-dimension structure and the structure of $M\setminus L$ as areas with extensive partial progress, and reports that Cusick's conjecture has been resolved, without stating that conjecture or citing a result. These references to progress are retained as supplied claims; they do not resolve the entire stated research direction.

The registry status is open and its research classification is OPEN-TRIAGE, whereas the summary describes a literature-survey program with many solved subquestions and no single complete resolution. No particular new subquestion is invented to replace the original direction.
\subsection{Short English statement}
Investigate the geometry of the Markov spectrum, viewed in its classical normalization as a subset of the real line. This is a broad research program, with reported progress on fractal dimensions and the difference from the Lagrange spectrum, rather than one conjecture with a specified completion criterion.
\subsection{Sources}
\begin{itemize}
\item Local UnsolvedMath registry, permanent ID 1900021, source code catalogue entry, ``Geometry of Continued Fractions --- Further open questions''; status open, classification OPEN-TRIAGE, updated 2026-08-06. Its research summary describes a literature survey with extensive partial progress and claims resolution of Cusick's conjecture, without a theorem statement or citation.
\item Karpenkov, \emph{Open problems in geometry of continued fractions} (2017), Problem 21, \url{https://arxiv.org/abs/1712.01450}. The record reports source-TeX extraction and access on 2026-07-29; it says the source presents an open item and current status requires release review. Rights: NEEDS\_REVIEW; public arXiv source TeX inspected according to the record.
\end{itemize}
\end{document}
